% Generated draft: requires mathematical and source review.
\documentclass[11pt]{article}
\usepackage[margin=25mm]{geometry}
\usepackage{fontspec}
\setmainfont{Latin Modern Roman}
\usepackage{amsmath,amssymb,mathtools,unicode-math}
\setmathfont{Latin Modern Math}
\usepackage{enumitem,needspace,xurl,hyperref}
\setlength{\parindent}{0pt}
\setlength{\parskip}{5pt}
\setlength{\emergencystretch}{3em}
\newcommand{\N}{\mathbb{N}}
\newcommand{\Z}{\mathbb{Z}}
\newcommand{\Q}{\mathbb{Q}}
\newcommand{\R}{\mathbb{R}}
\newcommand{\C}{\mathbb{C}}
\newcommand{\F}{\mathbb{F}}
\newcommand{\A}{\mathbb{A}}
\newcommand{\PP}{\mathbb{P}}
\newcommand{\E}{\mathbb{E}}
\newcommand{\eps}{\varepsilon}
\newcommand{\dd}{\,\mathrm{d}}
\begin{document}
\section*{UnsolvedMath ID: 20003113}
\textbf{Source problem code:} catalogue entry
\subsection{Definitions and mathematical statement}
Let $\Gamma<PSL(2,\R)$ be an infinitely generated Fuchsian group of the first kind, meaning that its limit set is the whole ideal circle. In the surface interpretation take a torsion-free group and the complete hyperbolic quotient $X=\mathbb H^2/\Gamma$ of infinite area, with no boundary. The original phrase ``closed hyperbolic surface of infinite area'' cannot mean compact closed; its intended meaning is boundaryless. Boundarylessness alone is not silently asserted equivalent to first kind for every infinite-area surface.

The source asks whether one can quasiconformally deform such a structure so that every essential simple closed curve is nonlengthened. For a marked deformation $f:X\to Y$, the requested inequality is
\[
 \ell_Y(f(\gamma))\leq\ell_X(\gamma)
 \qquad\text{for every essential simple closed curve class }\gamma.
\]
A nontrivial deformation must be intended, since the identity satisfies the inequalities; the source does not state whether this asks for one example, for every base surface, for a whole path, or merely another marked point. It mentions recent results of Alessandrini and Gu\'eritaud without a citation.

The supplied aid gives an obstruction for a specific natural family, on any complete hyperbolic surface where its curve hypotheses hold. Let $a$ be an essential simple closed geodesic with a simple transversal $b$, so $i(a,b)>0$. Let $X_t$ be a pure Fenchel--Nielsen twist along $a$, with $X_0=X$ and all other data unchanged. For every nonzero real twist parameter $t$, it reports infinitely many simple curve classes lengthened and infinitely many shortened relative to $X$.

More precisely, let $D_a$ be a Dehn twist and $b_n=D_a^{-n}(b)$ for $n\in\Z$. The reported tail certificate says that for each $t\ne0$ there is $N$ such that, after possibly interchanging the positive and negative tails according to the twist-sign convention,
\[
 \ell_{X_t}(b_n)>\ell_X(b_n)\quad(n\geq N),\qquad
 \ell_{X_t}(b_n)<\ell_X(b_n)\quad(n\leq-N).
\]
The threshold may depend on $X,a,b,t$. The aid also reports the two-sign infinitesimal analogue for the length derivatives of these twist-orbit curves at $t=0$, without providing its estimates. Thus a nonzero pure twist is not an all-nonincreasing deformation, even though it is a standard quasiconformal deformation.

In the one-sided convention
\[
 K(X,Y)=\log\sup_{\gamma}
 \frac{\ell_Y(\gamma)}{\ell_X(\gamma)},
\]
the lengthening tail implies $K(X,X_t)>0$ for every $t\ne0$. The shortening tail prevents the reversed comparison from satisfying all nonincrease inequalities as well. This is a pure-twist obstruction with an explicit transversal assumption; it does not rule out more general coordinated quasiconformal deformations or solve the source's full first-kind question. The dotted separator is layout, not a condition.
\subsection{Short English statement}
The source concerns nontrivial quasiconformal shortening on a complete infinite-area surface, with its quantifiers left partly open. The aid excludes pure twists: opposite tails of a Dehn-twist orbit supply lengthening and shortening curves. This special obstruction does not exclude every possible deformation.
\subsection{Sources}
\begin{itemize}
\item Local UnsolvedMath registry, permanent ID 20003113, source code catalogue entry; \texttt{aim-topology-notes.json}, notes[200].
\item AIM, \emph{Lipschitz metric on Teichmueller space}, Problem 4.1: \url{https://aimath.org/pastworkshops/lipschitzteichproblems.pdf}.
\item Original source mentions recent results of Alessandrini and Gu\'eritaud, without titles, dates or theorem statements. No reference for the pure-twist tail theorem is supplied by the research summary.
\item Machine-generated research attempt 1: partial result, open at the research run, rigor labelled ``proved,'' low novelty confidence; registry partially solved. Not peer reviewed; independently verify before citation. Registry updated 14 August 2026.
\end{itemize}
\end{document}
